%%%%%%%%%%  scan idx 19  |  folio 6  %%%%%%%%%%
\subsection*{3.}
Montrons maintenant comment on peut encore donner un sens à la formule (2.2) lorsqu'on abandonne l'hypothèse de régularité faite sur $X,Y$, en la remplaçant par une hypothèse de nature locale sur le morphisme $f$, que nous allons préciser.

\subsubsection*{3.1}
Une première difficulté provient du fait que l'observation qui avait servie de point de départ à la théorie [RRR], savoir que le groupe $K(X)$ des classes de faisceaux sur $X$ est le même, qu'on le définisse en termes de faisceaux cohérents, ou en termes de faisceaux localement libres, est liée de façon essentielle à l'hypothèse de régularité sur $X$. Dans le cas où $X$ est un schéma (disons noethérien, pour simplifier) général, il convient de distinguer entre ces deux groupes, que nous noterons respectivement $K_\bullet(X)$ et $K^\bullet(X)$, le premier ayant un caractère covariant en $X$ pour les morphismes propres, le deuxième un caractère contravariant en $X$ pour les morphismes quelconques. On notera que $K^\bullet(X)$ est muni de façon naturelle d'une structure d'anneau (et même de $\lambda$-anneau augmenté), alors que sur $K_\bullet(X)$ il n'y a plus en général une structure d'anneau, mais seulement une structure de module sur $K^\bullet(X)$. L'homomorphisme $x\to x\,\cl_\bullet(\underline O_X)$ (où $\cl_\bullet(\underline F)$ désigne la classe de $\underline F$ dans $K_\bullet(X)$) de $K^\bullet(X)$ dans $K_\bullet(X)$ n'est plus en général un isomorphisme. \uline{Pour notre formulation du théorème de Riemann--Roch généralisé, nous travaillerons exclusivement avec $K^\bullet(X)$, $K^\bullet(Y)$, à l'exclusion de $K_\bullet(X)$, $K_\bullet(Y)$}. Comme nous avons signalé dans 2.3., ces $\lambda$-anneaux augmentés, munis de leur $\lambda$-filtration naturelle, donnent naissance à des anneaux gradués $\Gr^\bullet(X)$, $\Gr^\bullet(Y)$, et des homomorphismes «~classes de Chern~»
\[
c_i:K^\bullet(X)\longrightarrow \Gr^\bullet(X),
\]
et de même pour $Y$. La difficulté vient cependant dans la définition d'un homomorphisme image directe
\begin{equation}\tag{3.1}
f_*:K^\bullet(X)\longrightarrow K^\bullet(Y),
\end{equation}
qui ne peut plus être définie par la formule naïve
\[
f_*(\cl^\bullet(\underline E))=\sum_i(-1)^i\cl^\bullet\bigl(R^if_*(\underline F)\bigr),
\]
car même si $\underline F$ est localement libre, les Modules $R^if_*(\underline F)$ ne sont plus en général localement libres, ni même de tor-dimension finie, et la formule précédente perd son sens. Une solution de cette difficulté est fournie par le point de vue des \underline{catégories dérivées}. En effet, alors que les $R^if_*(\underline F)$ sont en général des «~mauvais~» faisceaux, de tor-dimension infinie disons,
